\documentclass[aps,preprint,nofootinbib,preprintnumbers,eqsecnum,superscriptaddress]{revtex4}

\usepackage{amsmath} 
\usepackage{graphicx} 
\usepackage{amsthm}
\usepackage{amssymb} 
\usepackage{dsfont}
\usepackage{yfonts}
\usepackage{hyperref}
%\usepackage{floatrow,dcolumn,rotating}
\usepackage{array,xcolor,graphicx}
\usepackage{booktabs,multirow}
\usepackage[utf8]{inputenc}
\usepackage{mathtools}

\usepackage{slashed}
\usepackage{caption}
\usepackage{stackrel}
\usepackage{cases}
\usepackage{tcolorbox}
%\usepackage{subcaption}

%ADDED BY ALESSIO
\usepackage{enumerate}
%END ADDITION BY ALESSIO

%Ensuring that header of section+subsection is on the right
\usepackage{etoolbox}
\patchcmd{\section}
  {\centering}
  {\raggedright}
  {}
  {}
\patchcmd{\subsection}
  {\centering}
  {\raggedright}
  {}
  {}
%

\usepackage[normalem]{ulem}
\usepackage{mathtools}


%package for drawing tool
\usepackage[all]{xy}
\usepackage{tikz}
\usetikzlibrary{arrows.meta}

%Set colour for the hyperlink
\hypersetup{colorlinks=true,linkcolor=blue,citecolor=red,urlcolor=blue, linktocpage=true}

%%%%%%%%%%%%%%%%%%%%%%%%%%%%%%%%%%%%%%%%%%%


\newcommand{\be}{\begin{equation}}
\newcommand{\ee}{\end{equation}}
\newcommand{\bea}{\begin{eqnarray}}
\newcommand{\eea}{\end{eqnarray}}
\newcommand{\half}{\frac{1}{2}}
\newcommand{\ehalf}{\frac{e}{2}}
\newcommand{\ihalf}{\frac{i}{2}}
\newcommand{\quart}{\frac{1}{4}}
\newcommand{\G}{\Gamma}
\renewcommand{\O}{\Omega}
\renewcommand{\S}{\Sigma}
\newcommand{\U}{\Upsilon}
\renewcommand{\a}{\alpha}
\renewcommand{\b}{\beta}
\newcommand{\e}{\epsilon}
\newcommand{\g}{\gamma}
\renewcommand{\l}{\lambda}
\newcommand{\s}{\sigma}
\newcommand{\tB}{\tilde{B}}
\newcommand{\tC}{\tilde{C}}
\newcommand{\tR}{\tilde{R}}
\newcommand{\tS}{\tilde{S}}
\newcommand{\tL}{\tilde{L}}
\newcommand{\ka}{\kappa}
\newcommand{\derw}{{\bf\nabla}_w}
\newcommand{\bA}{\overline{A}}
\newcommand{\bC}{\overline{C}}



\newcommand{\nn}{\nonumber\\}
\newcommand{\bln}{\begin{align}}
\newcommand{\eln}{\end{align}}
\newcommand{\bst}{\begin{split}}
\newcommand{\est}{\end{split}}
\newcommand{\bi}{\begin{itemize}}
\newcommand{\ei}{\end{itemize}}
\newcommand{\ben}{\begin{enumerate}}
\newcommand{\een}{\end{enumerate}}

\def\AA{{\mathcal A}}
\newcommand{\BB}{{\mathcal B}}


\def\lad{L} %rob {l_{AdS}} changed notation because this is not the ads scale
\def\lg{\lam_{GB}}
\def\ns{N_{\sharp}}
\def\tf{\tilde{f}}
\def\w{\tilde{\om}}
\def\q{\tilde{q}}

\def\ov{\over}
\def\le{\left}
\def\ri{\right}
\def\ha{{1\over 2}}

\def\al{{\alpha}}
\def\ket#1{|#1\rangle}
\def\bra#1{\langle#1|}
\def\vev#1{\langle#1\rangle}
\def\det{{\rm det}}
\def\tr{{\rm tr}}
\def\Tr{\mathrm{Tr}}
\def\NN{{\cal N}}
\def\th{{\theta}}




\def \ra {\rightarrow}


\def\ep{{\epsilon}}
\def\apr{{\alpha'}}
\newcommand{\p}{\partial}
\def\LL{{\cal L}}
\def\TT{{\cal T}}
\def\tir{{\tilde r}}

\newcommand\ga{{\ensuremath{{\gamma}}}}
\newcommand\Ga{{\ensuremath{{\Gamma}}}}


\newcommand\sig{\sigma}
\newcommand\Sig{\Sigma}
\newcommand\lam{\lambda}
\newcommand\Lam{\Lambda}
\newcommand\om{\omega}
\newcommand\Om{\Omega}
\newcommand\vt{\vartheta}
\newcommand\de{{\ensuremath{{\delta}}}}
\newcommand\De{{\ensuremath{{\Delta}}}}
\newcommand\vp{\varphi}

\def\GeV{{\rm GeV}}
\def\fm{{\rm fm}}
\def\MeV{{\rm MeV}}
\def\lam{{\lambda}}

%\def\abs#1{\left| #1\right|}
\def\bR{\mathop{\bf R}}
%\def\bra#1{\left\langle #1\right|}
\def\eeq{\end{equation}}
\def\half{\frac{1}{2}}
\def\Im{\mathop{\rm Im}}
%\def\ket#1{\left| #1\right\rangle}
\def\Re{\mathop{\rm Re}}
\def\Tr{\mathop{\rm Tr}}
\def\crbig{\\\noalign{\vspace {3mm}}}
\newcommand{\HHH}{{\cal H}}

\def\tg{{\tilde g}}
\def\tR{{\tilde R}}
\def\tnab{{\tilde \nabla}}
\def\Rb{{[R^B]}}
\def\htt{{\hat t}}
\def\hx{{\hat x}}
\def\hy{{\hat y}}
\def\hz{{\hat z}}
\def\hr{{\hat r}}
\def\DD{{\cal D}}
\newcommand{\reef}[1]{(\ref{#1})}
\newcommand{\ssc}{\scriptscriptstyle}
\newcommand{\eg}{{\it e.g.,}\ }
\newcommand{\ie}{{\it i.e.,}\ }
%\newcommand{\comment}[1]{{\bf [[[#1]]]}}
\newcommand{\X}[5]{X^{\ssc (#1)}_{\ #2 #3}{}^{#4 #5}}

\newcommand\sA{{\ensuremath{{\mathcal A}}}}
\newcommand\sB{{\ensuremath{{\mathcal B}}}}
\newcommand\sC{{\ensuremath{{\mathcal C}}}}
\newcommand\sD{{\ensuremath{{\mathcal D}}}}
\newcommand\sF{{\ensuremath{{\mathcal F}}}}
\newcommand\sI{{\ensuremath{{\mathcal I}}}}
\newcommand\sG{{\ensuremath{{\mathcal G}}}}
\newcommand\sH{{\ensuremath{{\mathcal H}}}}
\newcommand\sL{{\ensuremath{{\mathcal L}}}}
\newcommand\sM{{\ensuremath{{\mathcal M}}}}
\newcommand\sN{{\ensuremath{{\mathcal N}}}}
\newcommand\sO{{\ensuremath{{\mathcal O}}}}
\newcommand\sV{{\mathcal V}}
\newcommand\sJ{{\mathcal J}}
\newcommand\sS{{\mathcal S}}
\newcommand\sR{{\ensuremath{{\mathcal R}}}}

\newcommand\tal{{\tilde \alpha}}
\newcommand\ta{{\tilde a}}
\newcommand\td{{\tilde d}}
\newcommand\tA{{\tilde A}}
\newcommand\tsD{{\tilde \sD}}
\newcommand\tG{{\tilde G}}
\newcommand\tow{{\tilde \om}}
\newcommand\tk{{\tilde k}}
\newcommand\tj{{\tilde j}}

\newcommand\bc{{\bar c}}
\newcommand\bb{{\bar b}}
\newcommand\bw{{\bar w}}
\newcommand\ba{{\overline a}}

\newcommand\bpsi{{\bar \psi}}
\newcommand\bPsi{{\overline \Psi}}
\newcommand\bPhi{{\overline \Phi}}

\newcommand\da{{\dagger}}
\newcommand\nab{{\nabla}}
\newcommand\Th{{\Theta}}
\def\th{{\theta}}

\newcommand\uz{{\underline{z}}}
\newcommand\ut{{\underline{t}}}
\newcommand\ur{{\underline{r}}}
\newcommand\ui{{\underline{i}}}
\newcommand\uj{{\underline{j}}}
\newcommand\umu{{\underline{\mu}}}
\newcommand\uM{{\underline{M}}}
\newcommand\utau{{\underline{\tau}}}
\newcommand\id{{\bf 1}}

\newcommand\bg{{\overline{g}}}
\newcommand\bG{{\overline{G}}}
\newcommand\bsig{{\overline{\sig}}}
\newcommand\bnab{{\overline{\nabla}}}
\newcommand\bSig{{\overline{\Sig}}}
\newcommand\bz{{\overline{z}}}

\numberwithin{equation}{section}

\begin{document}

\title{Non-invertible symmetries to study 3d complex scalar boson}
\author{Alessio Martini} 
\email{a.martini@uva.nl}
\affiliation{Institute for Theoretical Physics, University of Amsterdam\\} 
\begin{abstract}
To be written.
\end{abstract} 
\maketitle


\tableofcontents

\newpage

\section{Introduction}

To be written. 

\subsection{Motivation: ...} \label{sec:motivation} 

To be written.

\newpage

\section*{- - - Sparse Notes - - -}

Here I am writing things that will appear in the final thesis. But since there is no conclusion yet, and a unique logic that connects my journey into the physics topics I have encountered, I will write down here single topics that will appear in the final thesis, without showing a connection between them.

\section{What is a Generalized Global Symmetry?} \label{sec:gensym_intro} 

Let's stop and think. 
\ben[A.]
    \item What is a symmetry?
    \item What is generalized?
    \item Can we have a non global symmetry?
\een

\subsection{Let's start from the basics, with a reformulation}
\label{sec:reformulation_symemtry}

The answer to the first question is: "it is a topological operator". Or slightly more precise: the existence of an (ordinary) SYMMETRY (in Euclidean signature) implies the existence of a ($d-1$-dimensional) TOPOLOGICAL OPERATOR, but to understand this we need to revise what we already know from classical field theory.

Given an ordinary global continuous internal symmetry with (Lie) group $G$ (and Lie algebra $\frak{g}$ with generators $\{e^a\}_{a\in \{1,..., dim G \}}$), the Nöether's theorem tells us that we have conserved local currents $j^a_\mu (x):\mathcal{M} \to \mathcal{N}$ \textcolor{red}{I'M NOT SURE OF N} (it will be clearer to think of the current as a $\frak{g}$-valued 1-form $j^a=j^a_\mu dx^\mu$ globally defined).
\textcolor{red}{IS IT GLOBALLY DEFINED THIS EASILY? MAYBE IS A SECTION OF A BUNDLE}

Notice that the conservation equation $g^{\mu \nu}\partial_\mu j_\nu^a = 0$ is equivalent to $d \star j^a=0$. \footnote{Since the Hodge operator $\star:\Omega^p(\mathcal{M}) \to \Omega^{d-p}(\mathcal{M})$ makes manifest the metric dependence, we decided to keep it manifest in the index notation.} Thus, given a continuous ordinary global symmetry, we have a closed $\frak{g}$-valued $d-1$-form.

As a $d-1$-form, it naturally integrates over $d-1$-dimensional submanifold $\mathcal{M}_{d-1} \subset \mathcal{M}$: $Q^a(\mathcal{M}_{d-1}) = \int_{\mathcal{M}_{d-1}}\star j^a $, and as a closed form the dependence on $\mathcal{M}_{d-1}$ is said to be "topological" in the sense that it is independent with respect to continuous deformation of it: $\int_{\mathcal{M}_{d-1}}\star j^a = \int_{\mathcal{M}_{d}} d \star j^a - \int_{\mathcal{M'}_{d-1}}\star j^a = - \int_{\mathcal{M'}_{d-1}}\star j^a$ (where $\partial\mathcal{M}_{d}=\mathcal{M}_{d-1} \cup \mathcal{M'}_{d-1}$). \footnote{The minus sign is given by the opposite induced orientations from $\mathcal{M}_{d}$ to the boundary submanifolds.} Where we assumed that $j^a$ is not sourced by anything in $\mathcal{M}_{d}$, then it satisfies the continuation equation $d\star j^a=0$.

The $Q^a$s play the role of generators of the Lie algebra $\frak{g}$. From Lie theory, we can recover the connected part of the Lie group $G$ that includes the identity element via the exponential map: $U_g(\mathcal{M}_{d-1}) = \exp \{ i R(g)_a Q^a(\mathcal{M}_{d-1}) \}$ \footnote{The $i$ in the exponent is old well known a physical convention. Assumed that, requiring unitarity of $U_g$ corresponds to requiring hermitianity of $Q^a$ (opposed to antihermitianity, without the $i$), which we like in quantum mechanics since means that the eigenvalues (the values the observable can assume) are real, opposed to purely imaginary, again a convention we like to mantain \textcolor{red}{IS THE LATTER A CONVENTION?}.} where $R$ is taken to be "a" representation of the algebra. \textcolor{red}{HERE I'VE CHOSEN A REPRESENTATION (UNITARY?), I'M NOT SURE IF IT IS NECESSARY, SHOULD NOT}

Now that we have introduced the symmetry operator, and we have noticed the first property, namely being topological, let's notice the second: the composition (or more generally the fusion) of two symmetry operators respects a group law. $\lim_{\mathcal{M'}\to\mathcal{M}}U_{g'}(\mathcal{M'})U_{g}(\mathcal{M}) = U_{g'g}(\mathcal{M})$. Notice that the group does not need to be Abelian.

\textcolor{red}{CHARGED OBJECTS}

Here, the structure of the abstract symmetry is encoded in: what are the objects (the topological operators) and how they fuse (the group structure). Meanwhile, the structure of the objects charged under this symmetry is given by all the possible representations of the symmetry group.

\textcolor{red}{SOME DIAGRAMS, CONSIDER THE RIGHT CATEGORY, $Vec_k (G)$ ? AND $Rep(G)$ ?}

\textcolor{red}{HIGHER CATEGORICAL STRUCTURE} p.54 \cite{schafernameki2023ictplecturesnoninvertiblegeneralized}, for $d \geq 4$ the possible $d-1$-dim TQFTs we can stack are much richer (e.g G-symmetric modular tensor categories, i.e. non trivial topological order) TWISTED THETA DEFECT p.60 \cite{schafernameki2023ictplecturesnoninvertiblegeneralized}.

\subsection{Let's reverse the logic!}
The answer to the second question is: "two things, the dimension of the topological operator is no more constrained to be $d-1$ and the composition law does not to be group like anymore".

These generalizations come from the inversion of the following logic: the traditional idea of a continuous symmetry can be encoded into a topological operator and the fusion rules it satisfies. That is, from now on, we will start from topological operators of every dimension and we will call it "a symmetry operator", we are immediately going to see that this leads to a much wider variety of structures compared to the traditional "ordinary global symmetry" point of view.

The first generalization is: given that we find a $d-p-1$-dimensional topological operator, we will say our theory has a $p$-form symmetry, where the largest charged objects that transform under this symmetry are $p$-dimensional.

\textcolor{red}{DRAWINGS}

The second generalization is: the object obtained by fusing to topological operators is no more required to be a "simple" topological operator, where simple means it cannot be expressed as a sum of other topological operators. \footnote{Notice that this makes sense because "the sum" is not invertible, it does not exist a subtraction. Then we can decompose a "non simple" topological operator as a sum of simple ones.} In formulas, this can be written as: $\lim_{\mathcal{M'}\to\mathcal{M}}U_{a}(\mathcal{M'})U_{b}(\mathcal{M}) = \sum_c N^c_{ab} U_{c}(\mathcal{M})$. Here, $U_{i}(\mathcal{M})$ are 'simple' topological operator of fixed dimension. We will expand a little on the concept of 'simplicity' within the paragraph \ref{sec:operator_vs_defect}.

\subsection{[NOT REALLY RELEVANT...] That is not possible! Local symmetry means gauge invariance, a redundancy}

\textcolor{blue}{HERE I WANT TO MENTION SYMMETRY=NON TRIVIAL CHARGE, HAMILTONIAN CONSTRAINS. CONTINUOUS LOCAL INVARIANCE MEANS THAT WE ARE DESCRIBING THE THEORY WITH MORE DEGREES OF FREEDOM THAT THE PHYSICAL ONES. 3D SCALAR BOSON - GAUGE THEORY DUALITY, LINEAR THEORY-KLEIN GORDON HAS INFINITE SYMMETRIES PARAMETRIZED BY FUNCTIONS, CONSERVATION OF FOURIER MODES} \cite{Barcel__2016}

\section{ASIDE: What is the difference between a topological operator and a symmetry defect?}\label{sec:operator_vs_defect}

\textcolor{blue}{TIME ORIENTATION ETC.}

\section{Anomaly, from what it is to what it can mean} \label{sec:anomaly} 

\subsection{[ANOMALY THEORY] A modern formalism to describe anomalies}
For continuous symmetries we can compute anomalies from triangle current diagrams, for discrete symmetries we need refinement: anomaly theory (and polynomial) p.25 \cite{schafernameki2023ictplecturesnoninvertiblegeneralized}.

\begin{tcolorbox}

{\bf General Idea}

Given a QFT $\mathcal{T}$, with a symmetry $\mathcal{S}$ given by a group $G$ and an anomaly ('t Hooft), we can couple the theory to a classical background gauge field $B$. Then the partition function is not invariant under background gauge transformations: $Z[B+d\lambda]=Z[B]\times \exp \{ \int_\mathcal{M}f(B,\lambda)\}$.

Note that given this feature, the theory does not admit gauging with respect to the symmetry $G$.



\end{tcolorbox} 

\subsection{[IR CONSTRAINT] If the theory has an anomaly, we cannot have a trivial low energy effective theory}

\subsection{[MIXED ANOMALY INTRO] The anomaly can be shared between multiple symmetries, mixed anomalies}

\subsection{What happens if we gauge part of it?}
\textcolor{blue}{NON INVERTIBLE DEFECTS FROM GAUGING MIXED ANOMALIES:}  \cite{Kaidi_2022, Bhardwaj_2023}. References taken from p.55 of \cite{schafernameki2023ictplecturesnoninvertiblegeneralized}.

To be written.

\subsection{[HIGHER GROUP OR NON-INVERTIBLE DEFECT] What other things can happen if we have a mixed anomaly?}
THE FOLLOW IS TAKEN FROM \cite{Tachikawa_2020} AND EXPLAINS WHAT HAPPENS IF YOU HAVE A THEORY $\G$-SYMMETRIC WITH FINITE NORMAL NON-ANOMALOUS SUBGROUP $A$ AND YOU GAUGE $A$: \footnote{When the subgroup we gauge is non-normal it is more complicated, we know the answer for $d=2$ (footnote 1 of \cite{Tachikawa_2020})}
\begin{tcolorbox}

{\bf General Idea in general space time dimension}

Depending on how the $\G$-symmetry is anomalous, you have three possibilities:
\ben[\itshape i)]
    \item when $\G$ is anomaly free you get $\G /A \times \hat{A}$ with mixed anomaly determined by the extension $0 \to A \to \G \to \G /A \to 0$
    \item anomalous in a "tame" way, you get an extension of $\G /A$ by $\hat{A}$, ??$0 \to \hat{A} \to \G \to \G /A \to 0$?? leading to a higher-group  \footnote{Higher-group structure appears also gauging continuos symmetry with suitable anomaly (footnote 2 of \cite{Tachikawa_2020})}
    \item anomalous in a wilder way, you get a non-invertible sym leading to a higher-category structure
\een


\end{tcolorbox} 


\subsubsection{[HIGHER GROUP] gauging mixed anomalies, take 1}
2-groups and mixed anomalies are related by gauging \cite{Tachikawa_2020}. Reference taken from p.33 of \cite{schafernameki2023ictplecturesnoninvertiblegeneralized}.
IDEA: if you have a 2-group structure (with no anomalies), i.e. having a 1-form and a 0-form symmetries which are inter-related, and gauge the 1-form sym, you get the dual-pontryagin symmetry wrt the 1-form sym (a d-p-2-form sym, p.24 \cite{schafernameki2023ictplecturesnoninvertiblegeneralized}) will have a mixed anomaly (?with the original non-gauged 0-form sym?).

\subsubsection{[NON-INVERTIBLE DEFECT] gauging mixed anomalies, take 2}
\label{sec:mix_anom_noninv_defect}
AGAIN (end of \ref{sec:reformulation_symemtry}), if we have mixed anomaly (in $d \geq 4$), then the richer structure of TQFTs we can stack leads to TWISTED THETA DEFECTS. CLAIM: if we have a p-form sym in a mixed anomaly, and its background field appears linearly in the anomaly theory (d+1-dim theory) [in 3+1d-theories], then ...

THE FOLLOWING IS TAKEN FROM \cite{Kaidi_2022} APPENDIX D \& E.
\begin{tcolorbox}

{\bf General Idea}

Given a symmetry structure with a 0-form sym $G^{(0)}$ and a higher form sym $F^{(p)}$ with a mixed anomaly $\mathcal{A}_{d+1}=\int A^{(0)} \wedge f(B^{(p)})$, where $A^{(0)}$ and $B^{(p)}$ are the classical background gauge fields for the 0 and p-form sym respectively, and $f:H^{p+1}(\mathcal{M},\frak{f}) \to H^{d}(\mathcal{M},\frak{f})$. \textcolor{red}{DOES IT MAKE SENSE? CHECK! WE NEED f TO GIVE A TQFT.} Then, the symmetry defect of the 0-form sym $D^{(0)}=\exp \{ \int_{\mathcal{M}_{d-1}} \star j^{(0)}\}$ is not invariant under gauge sym transformation of the background $B^{(p)}$.

In fact, $\langle D(\mathcal{M_{d-1}}) \rangle = Z[A=\delta(\mathcal{M_{d-1}})]$...???... \textcolor{red}{CAN BE PHRASED LIKE THIS?}

Only $ D(\mathcal{M_{d-1}}) \times \exp \{ \int_{\mathcal{M}_d} f(B,\lambda)\}$ is gauge invariant. \textcolor{red}{TO BE CONCLUDED}


\end{tcolorbox} 

\begin{tcolorbox}

{\bf 1+1 dimension}

TOBEWRITTEN

\end{tcolorbox} 

\begin{tcolorbox}

{\bf 2+1 dimension}

TOBEWRITTEN

\end{tcolorbox} 

\begin{tcolorbox}

{\bf 3+1 dimension}

TOBEWRITTEN

\end{tcolorbox} 

\begin{tcolorbox}

{\bf ??4+1 dimension??}

new!

\end{tcolorbox} 

\section{[APPLICATION]}

\subsection{ABJ anomaly 4d}
This is an application of \ref{sec:mix_anom_noninv_defect}

\subsection{Complex scalar boson [3d only?]}
This theory has two phases, unbroken U(1) phase rotation, and superfluid: SSB U(1) where it does emerge/accident a U(1) (d-2)-form sym. Here in 3d, gauging the 0-form sym flips gives a gauge theory (?), flips the phase diagram (duality with inverse theory parameter?). 
GOAL: show that the (d-2)-form sym was already there in the unbroken phase under disguise (?higher group or non invertible defect?).
\textcolor{red}{WE MIGHT NEED TO INCREASE THE DIMENSIONS OTHERWISE THE TQFTS WE CAN STACK ARE SIMPLE, NO TWISTED THETA DEFECT} p.60 \cite{schafernameki2023ictplecturesnoninvertiblegeneralized}.

\newpage
 








\newpage
{\bf Acknowledgments:} I want to thank ... 

\bibliographystyle{utphys}
\bibliography{main}
\end{document}


\begin{tcolorbox}
{\bf Exercise:} 
Prove the adjoint trace identity (valid for any $SU(N)$ element):
\be
\Tr_{\mathrm{Adj}}(U) = \Tr_{f} U \Tr_{f} U^{\dagger} - 1 \label{adjtrace}  
\ee
Hint: use the fact that the adjoint representation can be obtained from a tensor product of fundamental and anti-fundamental reps: 
\be
f \otimes \bar{f} = \mathrm{Adj} \oplus {\mathbf{1}}
\ee
From \eqref{adjtrace} conclude explicitly that $\Tr_{\mathrm{Adj}}(g) = \Tr_{\mathrm{Adj}}(\mathbf{1})$ if $g \in \mathbb{Z}_N$. 
\end{tcolorbox} 

